\section{Results}

\subsection*{Cyclic topology in erlotinib-treated PC9 cells}

We first asked whether drug-treated NSCLC cells exhibit cyclic
topological structure. We analysed the PC9 erlotinib time-series
(GSE134839 \citep{aissa2021}; six timepoints D0--D11;
\cref{fig:study_overview}; Methods). All drug-treated timepoints
showed non-zero $H_1$ persistence, with maximum values ranging from
1.08 (D11) to 1.77 (D9) (\cref{fig:discovery_ph},
\cref{tab:per_timepoint}). Pairwise Wasserstein distances between
persistence diagrams separated D0 from all treated timepoints,
indicating that the distribution of topological features differs most
strongly between untreated and treated conditions
(\cref{fig:discovery_ph}). Per-timepoint sample sizes (200--379 cells)
limited individual permutation-test power; the larger validation
cohorts below recover significance (\cref{tab:validation}).

\subsection*{Osimertinib drives a monotonic topological increase across
cell line and patient-derived xenografts}

To test whether the cyclic signal strengthens with drug exposure, we
analysed a PC9 osimertinib time-series with Watermelon lineage tracing
(GSE150949 \citep{oren2021}; D0/D3/D7/D14; Methods). Max $H_1$ rose
monotonically from 1.41 at D0 to 3.78 at D14 (2.7-fold;
\cref{fig:validation}), with permutation tests confirming
$p < 0.01$ at D7 and D14 (\cref{tab:validation}).

To determine whether this finding extends beyond cell culture, we
analysed two EGFR-mutant PDX models (YU-003, YU-006) under matched
untreated and osimertinib-residual conditions (GSE243562
\citep{wurtz2024}; Methods). Both models showed increased max $H_1$
under residual disease: YU-003 from 2.79 to 3.85 (1.4-fold,
$p < 0.01$); YU-006 from 2.81 to 6.08 (2.2-fold, $p < 0.01$)
(\cref{fig:validation}, \cref{tab:validation}). This replicates the
cell-line finding in an \emph{in vivo} system with stromal, vascular,
and host-immune context.

A timepoint-label permutation test on osimertinib D14 vs.\ D0 further
confirmed that the observed topological change depends on drug exposure
rather than random point-cloud structure ($p < 0.02$; Methods).

\subsection*{The topological signal is not explained by cell-cycle
variation or batch effects}

Cell-cycle oscillation produces a cyclic snapshot signature that could
confound the topological signal. We applied two independent controls.
First, we removed cell-cycle genes from the feature set and recomputed
persistence; the ratio of max $H_1$ after to before removal ranged from
0.71 to 1.11 across all six timepoints
(\cref{fig:discovery_ph}, \cref{tab:cc_controls}). Second, we
regressed out S-phase and G2/M-phase scores from the expression matrix;
this \emph{increased} max $H_1$ in all five tested cohorts
(post/pre ratio 1.07--4.05; \cref{tab:cc_controls}).
Together, these results indicate that the observed cyclic topology is
not driven by cell-cycle variation (Methods).

To rule out batch or platform effects, we recomputed max $H_1$ on
Harmony-corrected embeddings \citep{korsunsky2019}. The osimertinib
monotonic trend, the YU-006 untreated-to-residual difference, and the
patient-tumour baseline all survived correction
(\cref{fig:harmony}; Supplementary Table~S10), indicating that the
topological signal reflects biology rather than technical artefacts.

\subsection*{Clinical validation in a 14-patient longitudinal cohort}

To establish a topological reference for human lung adenocarcinoma, we
analysed a treatment-naive patient LUAD atlas (GSE131907
\citep{kim2020}; Methods). Treatment-naive patient tumours showed
max $H_1 = 4.77$, substantially higher than treatment-naive cell line
(1.41) or PDX (2.79, 2.81) baselines (\cref{fig:master}). Drug
treatment moves cell-line and PDX topology from the baseline regime
toward the patient-tumour range (\cref{fig:master}; the within-system
monotonic increase is the primary claim, as cross-system comparisons
are confounded by platform and cell-count differences).

To directly test the framework on clinical tissue, we analysed the
EGFR-mutant subset of \citet{maynard2020} (14 patients with biopsies
at treatment-naive, persister, and progressive disease stages;
Methods). Pooled max $H_1$ grew monotonically across the three stages:
5.37 (treatment-naive) $\to$ 7.13 (persister) $\to$ 8.05 (progressive
disease) (1.5-fold; \cref{fig:patient}a). Three patients had matched
pre/post-treatment biopsies; max $H_1$ increased in all three
(median 1.6-fold, range 1.35--3.25-fold; Supplementary Table~S11),
supporting the use of max $H_1$ as a within-patient relative-change
readout.

\subsection*{EMT programs underlie cyclic topology}

To identify the transcriptional programs underlying loop formation, we
constructed Mapper graphs and computed gene-level connectivity scores
(Methods). In the cell-line analysis, top-connected genes were canonical
EMT markers (\emph{TPM1}, \emph{KRT8}, \emph{KRT19}); MSigDB Hallmark
EMT ranked first among enriched pathways. The PDX analysis recovered a
broader mesenchymal/secretory signature (\emph{S100A4},
\emph{LGALS3}, \emph{COL1A1}), consistent with an EMT-like plastic
state realised through a partially shared repertoire. Eleven genes
overlap between the two top-100 lists (hypergeometric
$p = 1.25 \times 10^{-11}$), indicating
significant but incomplete overlap between cell-line and PDX
plasticity programs. Differential expression recovered EMT-related
signatures in both systems, with system-specific secondary pathways
(Supplementary Tables~S2, S2$'$, and S4).

Using an unbiased PC1 filter in the Mapper, cell-cycle markers
co-ranked with EMT markers, indicating that the loops reflect a
multi-program state space rather than a pure EMT signature
(Supplementary Table~S2$'$). Variance-normalised connectivity scores
strengthened the EMT enrichment, confirming it is not an
expression-magnitude artefact (Supplementary Table~S9).

\subsection*{A generalisable, reproducible framework}

The pipeline comprises fifteen independent modules coordinated by a
single command (Methods). It ran without modification on Drop-seq and
10x cell-line data, 10x PDX data, and a 208{,}506-cell patient atlas.
Applying it to a new dataset requires only a processed expression
matrix with sample annotations and a choice of Mapper filter (PC1, EMT
score, or diffusion component; Methods). Resampling stability,
replication under alternative persistence summaries, and Mapper
parameter robustness are reported in Supplementary Methods~S10--S12.
